\subsection{Heights on the marked flag model}
\label{subsec:flag-point-heights}

The behavior near a wall will be inferred from heights on the compactified flag model. The height bounds must include the boundary as well as the open abelian locus.

Fix the deep cover at a nonterminal letter of either run, with CM
field \(E\), real degree \(f\), rank \(N\), and signature \((j,N-j)\).
Write \(A\) for its auxiliary ample metrized line and \(h_A\) for
the associated absolute point height. The perturbed line \(M_B(v,s)\)
has a summand \(sA\). Its root part has weight
\[
 w=v-sd_0,
\]
where
\(d_0\) is fixed on the cover. The arithmetic degree of a lattice over
a number field will always be divided by the degree of the field.
In this subsection constants may depend on the fixed cover. We shall
use a threshold in \(s\) independent of \(B\), of the point and its
field, and of the weight in a bounded set. This order of quantifiers
will allow us to take a limit for each cover before applying a
separate uniform estimate as the covers vary.

\paragraph{Lie lattices.}

The first height contribution is the degree of a line in the integral Lie algebra.

Let \(\mathcal L\) be the integral Lie algebra of a semiabelian model
over a number ring with principally polarized abelian generic fibre of
dimension \(g\). Use the polarization metric \(\pi H_\sigma\) on the
Lie algebra at each embedding, where \(H_\sigma\) is the Hermitian
form of complex uniformization. Put
\[
 h_*=-\widehat{\deg}\mathcal L
\]

and \(H=\max(h_*,0)\). For \(g\le g_0\), every integral sublattice
\(U\subset\mathcal L\) satisfies
\begin{equation}\label{eq:flag-lie-slope}
  \widehat{\deg} U\le \rk U\,C(g_0)\log(2+H).
\end{equation}
To see the logarithmic dependence, take a nonzero algebraic tangent
direction \(x\) and let \(a\) be the degree of the saturated line it
spans. Saturation of this line commutes with finite extension of the
number field, by flatness and the torsion-free quotient. Both its
normalized degree and its metric therefore persist over the fields
used below.

For an integer \(n\ge2\), put \(m=4n\) and let \(L\) be the
\(m\)-th power of the symmetric theta line, with its rigidification.
Over a finite extension take the semistable Moret--Bailly model
of \((A,L)\). The line has a cubical extension and the polarization
kernel of \(L^{\otimes 2}\) extends as a finite flat subgroup.
On this model Bost's semistable direct-image theorem gives a semistable lattice of sections with the cubical probability \(L^2\) metric. Its rank is 
\[
 \operatorname{rk}H^0(A,L)=m^g=(4n)^g.
\]
 In the Faltings normalization
of Gaudron's 2019 statement its slope is
\[
 -h_F/2+(g/4)\log(m/(2\pi^2)).
\]
For a principal polarization,
\[
 h_* = h_F+(g/2)\log\pi,
\]
because the tangent metric here is
\(\pi H_\sigma\). Thus the slope in these norms is
\[
 \nu=-h_*/2+(g/4)\log(m/(2\pi)).
\]

\cite{bost-mb-flag}\label{cor:flag-point-heights-1}

\begin{proposition}[Point-height bounds on the selected flag]
\label{prop:selected-flag-point-heights}
For the marked model the relative Lie lattices satisfy the
slope bound \eqref{eq:flag-lie-slope}. The resulting interior
and boundary point-height bounds use the same root metric.
On every component there is a dense set on which the two inequalities hold simultaneously
\[
 |h_A|\le C,\qquad |h_0(w)|\le C\|w\|_\infty\quad\text{for all }w.
\]
 The set is fixed before
\(B,v,s\).
\end{proposition}

\begin{proof}
\proofstep{Filtering by the evaluations at successively chosen points.}
Filter the space of sections by vanishing of the successive
directional derivatives at zero. The first \(n+1\) quotients have
rank one. Indeed \(\Theta^4\) separates the tangent at zero. Take two
sections of \(\Theta^4\), one nonzero at zero and the other with order
one along \(x\). Call the two sections \(s_0,s_1\). The powers satisfy
\[
 s_0^{\,n-i}s_1^i\in H^0(A,\Theta^{4n}),\qquad
          \operatorname{ord}_{x,0}(s_0^{\,n-i}s_1^i)=i,
                  \quad 0\le i\le n .
\]
Here the order is along the chosen tangent \(x\) at zero.
They exhibit a nonzero class in each of the successive
one-dimensional jet quotients.
Evaluate the \(i\)-th derivative on the kernel of the preceding
jets using the rigidification.

The derivative is integral on an integral tangent and an integral section. In the group law this derivative is the
coefficient of \(e_1\cdots e_i\) at the sum of \(i\) copies of that
tangent, in the algebra with \(e_k^2=0\). A change of trivialization
uses only its value at zero on this kernel. Writing \(j_i(s;x)\) for this derivative of a section \(s\) in the tangent \(x\), homogeneity of the multilinear coefficient gives
\[
 j_i(s;bx)=b^i j_i(s;x).
\]
Thus its finite-place image
pairs integrally with the \(i\)-th power of the saturated tangent
line.

At an embedding let \(\rho\) be the shortest nonzero period in
the \(H_\sigma\)-norm. The archimedean operator bound for this jet,
after removing \(\|x\|_{\pi H_\sigma}^i\), is
\[
  \beta_i(\rho)=(e m i)^{i/2}\sqrt{e\,g!}\,
       \max\{(2/(\sqrt\pi\rho))^g,m^{g/2}\}.
\]
Use \(1\) for the first factor when \(i=0\). In a translated
Appell--Humbert frame the centred Gaussian is
\(\exp(-\pi m H_\sigma(z,z))\) for the squared norm. A ball of
radius
\[
 \min(\rho/2,1/\sqrt{\pi m})
\]
injects up to its boundary
into the torus, and the Gaussian on it is at least \(e^{-1}\).

The holomorphic mean-value inequality and the volume of this ball
give the last two factors in the comparison of the supremum and
probability \(L^2\) norms. The fundamental period parallelepiped
has volume one in these coordinates by principality. Cauchy's
inequality along the unit tangent with radius
\(\sqrt{i/(\pi m)}\) in the \(H_\sigma\)-norm, and \(i!\le i^i\),
give the first factor with the stated tangent metric. The calculation
applies to every complex linear combination of sections in the
kernel. \label{cor:flag-point-heights-2}

\proofstep{The quotient heights and the product formula.}
The determinant of a rank-one quotient and the product formula
now give its normalized degree at most
\[
 -ia+\operatorname{mean}_\sigma\log\beta_i(\rho_\sigma).
\]

The remaining kernel has degree at most its rank times \(\nu\) by
semistability. Subtracting this from the full degree gives
\[
 (n+1)\nu\le -\tfrac12n(n+1)a+
           \sum_{i=0}^n\operatorname{mean}_\sigma\log\beta_i(\rho_\sigma).
\]
For principal polarization the matrix lemma bounds the mean of
\(\rho_\sigma^{-2}\) by \(14\max(1,h_*-(g/2)\log\pi,\log(g!))\).
Apply concavity of \(\log\) to \(1+\rho_\sigma^{-2}\). The mean in
the last display is at most
\[
  \tfrac i2\log(e m n)+\tfrac g2\log m+
        \tfrac g2\log(2+H)+O_g(1).
\]
Field extension repeats each base embedding equally often in this
mean.

Take \(n=2\lfloor H\rfloor+2\). Then \(H/n\le1\) and
\(\log n\le\log(2+H)+\log 2\). Dividing by \(n(n+1)/2\) gives
\(a\le\log(2+H)+O_g(1)\). Let \(U\) be a Hermitian sublattice of positive rank. The absolute
Siegel lemma gives, after finite extension, a line \(L\subset U\)
with
\[
    \widehat{\deg} L\ge \frac{\widehat{\deg} U}{\operatorname{rk}U}
                         -\tfrac12\log\operatorname{rk}U-1.
\]
The inequality extends to \(L\) from its saturation \(L^{\rm sat}\) by
\[
 \widehat{\deg} L\le\widehat{\deg} L^{\rm sat},\qquad \operatorname{rk}L=\operatorname{rk}L^{\rm sat}.
\]
This proves \eqref{eq:flag-lie-slope}, taking
a maximum over the finitely many dimensions up to \(g_0\). Rank zero
contributes zero. \cite{matrix-flag,absolute-siegel-flag}

\paragraph{Interior points.}
Take the semiabelian integral model of the abelian variety at the
given algebraic point after semistable extension. Its dimension is
\(fN\). Pull the marked flag back with the same point, intersect the
rational steps with the Lie lattice and give them the induced
metrics. If \(W_k\) denotes this intersection, the determinant comparison reads
\[
 h_{\rm step}=-\widehat{\deg}(\det W_k)+O(1).
\]
Indeed the product of the successive local quotient generators is
a determinant generator, and the product of their quotient norms is
the Gram determinant norm. At the finitely many primes where
integral eigenprojections require a denominator the field
discriminant bounds the change in the normalized valuation.

At infinity the polarized metric pulled back at the point
is the Hodge Gram form in its marked frame. The successive Gram
determinant ratios therefore agree at every embedding. If a rational determinant frame over \(K\) is rescaled by \(a\in K^\times\), the changes cancel by
\[
 \sum_v[K_v:\mathbb Q_v]\log |a|_v=0,
\]
with absolute values extending the standard ones. Applying the determinant comparison to the full Lie lattice gives
\(h_F=h_*+O(1)\). If \(h_{\rm step}\) denotes any one of the step
heights, \eqref{eq:flag-lie-slope} consequently gives
\[
                    h_{\rm step}\ge-C\log(2+\max(h_F,0))-C .
\]
Applying it to the entire lattice bounds the negative part of
\(h_F\) by the same quantity in absolute value.
\label{cor:flag-point-heights-3}

\proofstep{Root heights in the interior metric.}
At \(\tau\) write \(R_{\tau i}\), \(0\le i<N\), for the
graded root heights in the Hodge metric and
\[
 P_{\tau k}= \sum_{i<k} R_{\tau i}.
\]
The first \(j\) entries are the
graded Lie heights at \(\tau\). The remaining entries are
the negatives of those at \(\bar\tau\) in reverse order.
To bound the full sum \(P_{\tau N}\) in absolute value, use that
the full determinant and its inverse have integral sections up
to fixed denominators, nonzero at all generic-fibre points,
on the determinant cover. The smooth reference metrics bound their
norms above. Their changes to the Hodge metric for this central
weight are bounded, since its potential extends smoothly. Thus the step height in the first block reads
\[
 h_{\rm step}(\tau,k)=-P_{\tau k}\qquad(k\le j).
\]

For \(k>j\) it is the step height of dimension \(N-k\) at
\(\bar\tau\) minus \(P_{\tau N}\). Both have the preceding lower
bound. Moreover the sum of the two unsigned-block Lie determinants
in the Hodge metric, with negative sign, reads \(h_F+O(1)\) after
summing over \(\tau\). The two block sums differ by \(P_{\tau N}\).
Consequently
\[
  \sum_\tau -P_{\tau j}=\tfrac12 h_F+O(1).
\]
The field and level are fixed in these \(O(1)\) terms.
\label{cor:flag-point-heights-4}

Let \(\delta_{\tau k}=w_{\tau,k-1}-w_{\tau k}\) for \(1\le k<N\).
For a dominant \(w\) of norm at most \(2\), and
\(0\le b\le\min_\tau\delta_{\tau j}\), summation by parts gives
\[
 h_{\rm H}(w)=
 \sum_\tau\Big(\sum_{k=1}^{N-1}\delta_{\tau k}(-P_{\tau k})
                      -w_{\tau,N-1}P_{\tau N}\Big)
 \ \ge\ \tfrac b2\max(h_F,0)-C\log(2+\max(h_F,0))-C .
\]
To get the inequality, reserve \(b\) at each middle gap. The
remaining coefficients are nonnegative and their sum at each
embedding is at most \(4\). This also shows why the possibly
unbounded individual middle heights must first be added with a
common coefficient.

\paragraph{Boundary points.}
Use the smooth reference metric at a boundary point and the actual
extension by a torus of its polarized abelian part. After extending
the residue number field one obtains its Lie lattice \(L\), a map
\(q:L\to L'\) to the lattice of the abelian part, and a frame
\(\epsilon\) of the torus. The rational map is onto and its integral
kernel is the span of \(\epsilon\). On these data there is a fixed
\(D\ge1\) such that
\[
   {\rm Gram}(\epsilon)\ \ge D^{-1}I,\qquad
   {\rm Gram}(qv)\ \le D\,{\rm Gram}(v)
\]
for each tuple \(v\), at every embedding.

For the first comparison
split the torus over a finite cover of a dense open of each
boundary stratum. Its cocharacter basis is saturated in the
integral Lie algebra by the integral torus inclusion. Its norm
comparison on the split stratum is the torus-frame bound. Apply
Noetherian induction to the closed complement, keeping the maximum
of the finitely many bounds. The second comparison comes from the
local Hermitian bound on the map to the polarized quotient and a
finite cover of the compact complex charts. Both comparisons are
stable after extension to the common field of definition of the
point and frame. \label{cor:flag-point-heights-5}

\proofstep{The Lie quotient at the boundary.}
For \(W\subset L\) put \(W_0=W\cap\ker q\). Its coordinate lattice
in the torus basis has degree at most zero for the standard
metric, by a nonzero integral minor and the product formula.
Consequently
\[
 \widehat{\deg}W_0\le(\rk W_0)\log D/2.
\]

Using the quotient metric on \(W/W_0\), the second Gram inequality
gives
\[
 \widehat{\deg}W\le \widehat{\deg}W_0+\widehat{\deg}(qW)
                            +\tfrac12\rk(qW)\log D .
\]
In particular
\[
 h_*'=-\widehat{\deg}L'\le -\widehat{\deg}L+O(1),
\]
by taking \(W=L\) and using
\(qL\subset L'\) of full rank. The smooth determinant height
\(h_\omega=-\widehat{\deg}L+O(1)\) satisfies
\[
 |h_\omega|\le C(1+|h_A|)
\]
by ample height comparison.
Also \(h_A\ge-C\), using a finite generating family in an ample
power. Applying \eqref{eq:flag-lie-slope} to \(qW\) gives the
lower bound \(-C\log(2+C+h_A)-C\) for each smooth step height,
using the determinant comparison for the saturated marked steps
of \(L\). The argument includes a zero-dimensional abelian part
by taking \(qW=0\).

The smooth root heights have the same signed reversal as the
interior heights and bounded full sum at each \(\tau\). For
\(k>j\) use the conjugate step of length \(N-k\), which contributes
the same logarithmic loss. Summation by parts for dominant \(w\)
of norm at most \(2\) therefore gives
\[
 h_0(w)\ge -C\log(2+C+h_A)-C .
\]
Here \(h_0(w)\) uses the empty finite-prime set and smooth metrics.
\label{cor:flag-point-heights-6}

\paragraph{Upper heights on a dense set.}

The lower height bounds apply to arbitrary points. To bound the intersection in the other direction we need upper bounds on points that are dense in each component.

Fix a component, embed it by a power of the ample line, and
choose an affine chart on which the projective coordinate ratios
\(z_1,\ldots,z_m\) are regular. Noether normalization gives a
finite map
\[
 \pi\colon U\longrightarrow\mathbb A^d_K,\qquad
 K[t_1,\ldots,t_d]\hookrightarrow K[U],\qquad
 K(U)/K(t_1,\ldots,t_d)\text{ finite}.
\]
For \(0\ne a\in K[U]\), its field norm is nonzero.
Clear denominators and restrict to the open set over which the
norm specializes as the determinant of multiplication by \(a\).
The tuples of roots of unity are Zariski dense in \(\mathbb A^d\),
so take one in this open set on which the norm is nonzero. At that
tuple the determinant is invertible; every point of its nonempty
fibre lies in \(D(a)\). The union of the fibres over torsion tuples
therefore meets every nonempty principal open of \(U\).

Each \(z=z_j\) satisfies a monic equation
\(z^s+\sum_{i<s}f_i(t)z^i=0\) with \(f_i\in K[t]\).
Write \(f_i=\sum_\alpha c_{i\alpha}t^\alpha\). For a torsion tuple
and an absolute value above \(v\) the monic root bound gives
\[
 \max(1,|z|_v)\le
 \begin{cases}
  1+\sum_{i,\alpha}|c_{i\alpha}|_v,& v\mid\infty,\\
  \max(1,\max_{i,\alpha}|c_{i\alpha}|_v),& v\nmid\infty .
 \end{cases}
\]
Let \(C_v\) be the maximum of these bounds over the finitely
many coordinate ratios. It equals one outside a fixed finite
set. With the normalized local weights \(n_v\) of \(K\), summing
over the places of any point field above \(v\) gives
\begin{equation}\label{eq:root-fibre-height}
       h([1:z_1:\cdots:z_m])\le\sum_v n_v\log C_v.
\end{equation}
Passing from the projective height in \eqref{eq:root-fibre-height}
to the fixed ample reference metric changes the bound by \(O(1)\).
Lift these closed points to marked algebraic points using the
component's point map. Its image contains every such closed point,
which is the density test on the marked component.

On the same set every smooth root height \(h_x\) is bounded in
absolute value. On the common root model, twist the root and its
dual by sufficiently high powers of the pulled-back ample line so
that the resulting sections span on the generic fibre. Clear fixed denominators.
For each point use a section nonzero there. Bounded norms of the
finite list give both \(\pm h_x\le C_x(|h_A|+1)\).
Thus one dense set per component satisfies
\(|h_A|\le C\) and
\[
 |h_0(w)|\le C\|w\|_\infty
\]
simultaneously for all \(w\). The same dense set can thus be used as \(B,v,s\) vary in the subsequent intersection bounds.
\label{cor:flag-point-heights-7}
\end{proof}
